\section{滤过集上的函数比较}

设 E 是一个被具有基 F 的滤子所滤过的集合 (Gen.~Top., I, p.~58)；在本章中我们将考虑这样的函数：它们的定义域是属于滤子基 F 的 E 的一个子集（该子集依赖于函数），并且它们的值或者取在实数域 R 中，或者更一般地取在一个赋值域上的赋范向量空间中 (Gen.~Top., IX, p.~170)。

在应用中，E 大多是实空间 \(R^{n}\) 的一个子空间，或者是扩充实直线 \(\overline{R}\) ，而 \(F\) 则是 E 的某个闭包点的邻域滤子在 E 上的迹，或者是 E 的相对紧子集的补集所成的滤子（“无穷远点的邻域”）。

一般说来，仅知道这样一个函数沿 \(\mathfrak{F}\) 趋于一个给定的极限，并不足以处理由这个函数构成的表达式的所有“沿 \(\mathfrak{F}\) 过渡到极限”的问题。

例如，当实变量 \(x\) 趋于 \(+\infty\) 时，三个函数 \(x\) 、\(x^2\) 与 \(\sqrt{x}\) 都趋于 \(+\infty\) ，但在诸表达式

\[
(x + 1) ^ {2} - x ^ {2}, \qquad (x + 1) - x, \quad \sqrt {x + 1} - \sqrt {x},
\]

中，第一个趋于 \(+\infty\) ，第二个趋于 1，第三个趋于 0。

因此，重要的不只是知道函数沿 \(\mathfrak{F}\) 的极限值（当这极限存在时），还要知道函数接近其极限的“方式”；换句话说，我们要对趋于同一极限的函数所成之集进行分类。
% semantic-fragments: legacy-input-seam
\relax{}
